% Compilar con xelatex (nunca pdflatex).
% Estructura para la 2026년 동계학술논문발표회 — 한국통계학회
% Diciembre 22-23, 2026 · Universidad Chung-Ang (Seúl)
\documentclass[11pt,a4paper]{article}

\usepackage{fontspec}
% Malgun Gothic cubre Unicode completo (coreano + latín); no requiere xeCJK
\setmainfont{Malgun Gothic}[Ligatures=TeX]
\setsansfont{Malgun Gothic}[Ligatures=TeX]
\setmonofont{Courier New}

\usepackage{amsmath,amssymb}
\usepackage{graphicx}
\usepackage{float}
\usepackage{enumitem}
\usepackage{booktabs,longtable,array,multirow}
\usepackage{xcolor}
\usepackage[margin=2.5cm]{geometry}
\usepackage[hidelinks]{hyperref}
\usepackage{xurl}

% Término técnico trilingüe: \term{español}{한국어}{english}
\newcommand{\term}[3]{\textbf{#1} (\textcolor{gray}{#2 / \textit{#3}})}
% Sección pendiente de completar
\newcommand{\todo}[1]{\textcolor{red}{\textbf{[TODO: #1]}}}

\newcolumntype{C}[1]{>{\centering\arraybackslash}p{#1}}
\setlength{\parskip}{0.5em}
\setlength{\parindent}{0pt}

% ─────────────────────────────────────────────────────────────────────────────
\begin{document}

% ── Encabezado estilo KSS ────────────────────────────────────────────────────
{\footnotesize\raggedleft
한국통계학회 2026년 동계학술논문발표회\\
2026년 12월 22--23일 · 중앙대학교(서울캠퍼스)\par}

\vspace{1.2em}

\begin{center}
{\LARGE\bfseries
한국 대중가요 가사의 어휘 다양성과 영어 혼용 비율의\\[0.3em]
통계적 분석: 2003--2020년 Solid--Liquid 구조 변화\par}

\vspace{0.5em}
{\large\itshape
Análisis estadístico de la diversidad léxica y el ratio inglés\\[0.2em]
en letras de K-pop: cambio de estructura Solid--Liquid 2003--2020\par}

\vspace{1.2em}
{\large
\todo{Nombre del autor}\textsuperscript{*} \quad
\todo{Nombre del asesor}\textsuperscript{**}\par}
\end{center}

\vspace{0.5em}
\hrule
\vspace{0.5em}
{\footnotesize
\textsuperscript{*} Autor principal, estudiante de maestría, \todo{Departamento}, \todo{Universidad} (\todo{email})\\
\textsuperscript{**} Autor de correspondencia, \todo{Cargo}, \todo{Facultad}, \todo{Universidad} (\todo{email})\par}
\vspace{0.5em}
\hrule

% ── Resumen (coreano) ────────────────────────────────────────────────────────
\section*{요약 (Resumen)}

이 연구는 2003년부터 2020년까지 한국 음악 차트 상위권에 오른 대중가요 가사를
대상으로, 어휘 다양성과 영어 혼용 비율의 통계적 변화를 분석한다.
\term{시계열 분석}{시계열 분석}{time series analysis},
\term{스피어만 상관분석}{스피어만 상관}{Spearman correlation}
및 \term{엔트로피}{엔트로피}{entropy}·\term{HHI}{HHI}{Herfindahl--Hirschman Index}
기반의 \term{Solid--Liquid 분류}{집중형--분산형 분류}{Solid--Liquid classification}를
활용하여 가사 어휘 구조의 집중형(Solid)·분산형(Liquid) 특성이 연도에 따라
어떻게 변화하는지 살펴본다.
또한 \term{XGBoost}{XGBoost}{XGBoost} 회귀 모형을 통해
영어 비율, 어휘 다양성, 반복률 등의 언어적 변수가 차트 순위를 설명하는
추가적인 통계적 유의성을 검증한다.
분석 결과, 영어 혼용 비율은 2003년 이후 구조적으로 증가하는 추세이며,
어휘 분산 구조(Liquid)는 특정 연도군에 집중되는 경향이 있음을 확인하였다.
이 연구는 한국어 텍스트 코퍼스를 활용한 시계열 언어 통계 분석의
새로운 프레임워크를 제시한다.

\textbf{주요어:} 대중가요 가사, 어휘 다양성, 영어 혼용, 엔트로피, HHI,
Solid--Liquid, 시계열 분석, XGBoost, 차트 순위

\vspace{0.8em}

% ── Abstract (inglés) ───────────────────────────────────────────────────────
\section*{Abstract}

This study conducts a statistical analysis of the lexical diversity and
English-mixing ratio in Korean popular music (K-pop) lyrics spanning from
2003 to 2020.
Using a self-constructed longitudinal corpus derived from Korean music chart
data, we examine temporal changes in vocabulary concentration through
entropy and Herfindahl--Hirschman Index (HHI) metrics, classifying each year
as exhibiting a concentrated (Solid) or dispersed (Liquid) lexical structure.
Additionally, Spearman correlation and XGBoost regression are applied to assess
the statistical relationship between linguistic variables --- English ratio,
lexical diversity, and repetition rate --- and chart ranking positions.
Results indicate a structurally increasing trend in English mixing since 2003,
and a Liquid lexical structure concentrated in specific year clusters.
This study proposes a new time-series linguistic framework for statistical
analysis of Korean text corpora.

\textbf{Keywords:} K-pop lyrics, lexical diversity, English mixing, entropy,
HHI, Solid--Liquid, time series, XGBoost, chart ranking

% ── I. Introducción ─────────────────────────────────────────────────────────
\section*{I. Introducción}

La música popular coreana (K-pop) ha experimentado una transformación lingüística
notable durante las dos últimas décadas.
Mientras que en los primeros años del siglo XXI las letras se componían
predominantemente en coreano, a partir de mediados de la década de 2000 comenzó
a observarse una incorporación creciente de vocabulario en inglés, un fenómeno
que coincide con la expansión internacional del Hallyu (한류)
(Kim \& Ryoo, 2007; Jung, 2014).

Sin embargo, esta transformación lingüística no ha sido estudiada de forma
cuantitativa y longitudinal desde una perspectiva estadística formal.
La mayor parte de la literatura existente aborda el fenómeno del code-mixing
en el K-pop desde enfoques cualitativos o sobre corpus reducidos y acotados
temporalmente (Moody \& Matsumoto, 2011; Lee, 2004).
Existe, por tanto, un vacío de investigación en el análisis estadístico de la
evolución temporal de la estructura léxica de las letras de K-pop a gran escala.

Este estudio busca cubrir dicho vacío mediante el análisis de un corpus
longitudinal propio de letras de canciones que alcanzaron las posiciones
superiores del ranking mensual en el período 2003--2020, un total de
\todo{N canciones} con \todo{N tokens}.
Se proponen tres preguntas de investigación:

\begin{enumerate}[leftmargin=1.8em]
    \item ¿Existe una tendencia estadísticamente significativa de aumento del
    ratio inglés en las letras de K-pop durante el período 2003--2020?
    \item ¿Cómo varía la estructura léxica (Solid vs.\ Liquid) a lo largo de
    los años analizados, medida a través de entropía y HHI?
    \item ¿En qué medida las variables lingüísticas (ratio inglés, diversidad
    léxica, tasa de repetición) se asocian estadísticamente con la posición
    en el ranking mensual?
\end{enumerate}

% ── II. Marco teórico ───────────────────────────────────────────────────────
\section*{II. Marco teórico}

\subsection*{2.1 Code-mixing y contacto lingüístico en la música popular}

El \term{code-mixing}{코드 혼용}{code-mixing} es el fenómeno por el cual
hablantes alternan entre dos o más lenguas dentro de un mismo texto o discurso
(Myers-Scotton, 1993).
En el contexto de la música popular, el uso de inglés en letras de otros idiomas
ha sido documentado en múltiples mercados --- japonés, español, sueco ---
como estrategia de posicionamiento y ampliación de audiencias (Moody \& Matsumoto, 2011).
En el caso del K-pop, Lee (2004) identificó patrones recurrentes de alternancia
coreano-inglés en letras de grupos de idol de la década de 1990, y estudios
posteriores han relacionado este fenómeno con la globalización activa del género
(Jung, 2014).

\subsection*{2.2 Diversidad léxica y su medición estadística}

La \term{diversidad léxica}{어휘 다양성}{lexical diversity} cuantifica la riqueza
del vocabulario empleado en un texto.
Las medidas clásicas como la \term{relación tipo-token}{타입-토큰 비율}{type-token ratio}
(TTR) presentan sensibilidad al tamaño del texto (Templin, 1957), por lo que en
este estudio se complementa con medidas de distribución:

\[
H = -\sum_{i=1}^{K} p_i \log p_i, \quad
H_{norm} = \frac{H}{\log K}, \quad
\mathit{HHI} = \sum_{i=1}^{K} p_i^2
\]

donde $K$ es el número de tipos léxicos distintos y $p_i$ la proporción
de cada tipo sobre el total de tokens.
Siguiendo la clasificación propuesta por Moon y Pyun (2025), definimos:

\[
\mathit{Liquid} \Longleftrightarrow
(H_{norm} \geq 0.90) \wedge (\mathit{HHI} \leq 0.23) \wedge (n_{active} \geq 4)
\]

siendo $n_{active}$ el número de categorías léxicas con peso $\geq 10\%$.
Cuando no se cumple alguna condición, la estructura se clasifica como Solid.

\subsection*{2.3 Correlación lingüística y posición en el ranking}

La relación entre características textuales y popularidad ha sido explorada
principalmente en contextos de análisis de sentimiento y recomendación de contenidos
(Schedl, 2014).
En el ámbito de la música popular, Napier y Shamir (2018) identificaron
asociaciones entre simplicidad léxica y mayor desempeño comercial.
Este estudio extiende esa línea hacia el mercado coreano y un horizonte
temporal de 17 años.

% ── III. Datos ──────────────────────────────────────────────────────────────
\section*{III. Datos y preprocesamiento}

\subsection*{3.1 Descripción del corpus}

El corpus se construyó a partir de los rankings mensuales de un sitio de
música en streaming coreano, recopilando las canciones posicionadas en el
top-100 mensual durante el período enero de 2003 -- diciembre de 2020,
lo que comprende 216 meses de datos.

\begin{table}[H]
\centering
\small
\begin{tabular}{lr}
\toprule
\textbf{Característica} & \textbf{Valor} \\
\midrule
Período                       & 2003.01 -- 2020.12 \\
Total de meses                & 216 \\
Canciones únicas (aprox.)     & \todo{N} \\
Total de tokens (aprox.)      & \todo{N} \\
Idioma dominante              & Coreano (con presencia de inglés) \\
Herramienta de análisis morfológico & Kiwipiepy (Kiwi) \\
\bottomrule
\end{tabular}
\caption{Descripción general del corpus}
\end{table}

Para el análisis morfológico se empleó \texttt{kiwipiepy} (Kiwi),
analizador de morfemas del coreano basado en modelos de lenguaje,
que permite identificar sustantivos (etiquetas NN*), verbos y adjetivos
de forma robusta incluso en textos coloquiales como letras de canciones.

\subsection*{3.2 Variables construidas}

\begin{table}[H]
\centering
\small
\begin{tabular}{lll}
\toprule
\textbf{Variable} & \textbf{Descripción} & \textbf{Tipo} \\
\midrule
\texttt{english\_ratio}   & Proporción de tokens en inglés sobre el total & Continua [0,1] \\
\texttt{entropy}          & Entropía de la distribución de sustantivos & Continua $\geq 0$ \\
\texttt{HHI}              & Índice Herfindahl-Hirschman léxico & Continua (0,1] \\
\texttt{H\_norm}          & Entropía normalizada ($H / \log K$) & Continua [0,1] \\
\texttt{n\_active}        & Categorías léxicas con peso $\geq 10\%$ & Discreta \\
\texttt{repetition\_rate} & Proporción de tokens repetidos en la canción & Continua [0,1] \\
\texttt{rank}             & Posición en el ranking mensual (1--100) & Discreta \\
\texttt{solid\_liquid}    & Clasificación binaria Solid / Liquid & Categórica \\
\bottomrule
\end{tabular}
\caption{Variables del análisis}
\end{table}

\subsection*{3.3 Preprocesamiento}

Se aplicó un preprocesamiento mínimo sobre las letras crudas:
eliminación de caracteres especiales y URLs,
normalización Unicode,
y separación de tokens coreanos e ingleses mediante expresiones regulares
(\texttt{[가-힣]+} para coreano, \texttt{[a-zA-Z]+} para inglés).
No se eliminaron palabras vacías del coreano con el fin de preservar
la estructura rítmica y de repetición característica de las letras.

% ── IV. Metodología ─────────────────────────────────────────────────────────
\section*{IV. Metodología}

\subsection*{4.1 Análisis de tendencia temporal del ratio inglés}

Para verificar si el ratio inglés presenta una tendencia estadísticamente
significativa durante el período analizado, se aplica la
\term{prueba de Mann-Kendall}{Mann-Kendall 검정}{Mann-Kendall test},
un test no paramétrico para tendencias monotónicas en series temporales
(Mann, 1945; Kendall, 1975).
La estadística de prueba $S$ se define como:

\[
S = \sum_{k=1}^{n-1}\sum_{j=k+1}^{n} \mathrm{sgn}(x_j - x_k)
\]

donde $\mathrm{sgn}(\cdot)$ es la función signo.
La hipótesis nula $H_0$ es la ausencia de tendencia monotónica.

\subsection*{4.2 Clasificación Solid--Liquid por año}

Para cada año $y$ del período analizado se calcula la distribución
de tokens léxicos en $K$ categorías semánticas derivadas mediante
análisis de frecuencias de sustantivos.
A partir de dicha distribución se obtienen $H_{norm}(y)$, $\mathit{HHI}(y)$
y $n_{active}(y)$, y se clasifica el año como Solid o Liquid
según el criterio de la sección 2.2.

\subsection*{4.3 Correlación ranking-variables lingüísticas}

La relación entre las variables lingüísticas y la posición en el ranking
se cuantifica mediante la
\term{correlación de Spearman}{스피어만 상관계수}{Spearman rank correlation}:

\[
r_s = 1 - \frac{6\sum d_i^2}{n(n^2-1)}
\]

donde $d_i$ es la diferencia de rangos entre las dos variables para la
observación $i$.
Se reportan los coeficientes $r_s$ y sus $p$-valores correspondientes
con corrección de Bonferroni para comparaciones múltiples.

\subsection*{4.4 Modelo XGBoost: variables lingüísticas como predictores del ranking}

Para evaluar el poder explicativo conjunto de las variables lingüísticas
se construyen dos modelos de regresión XGBoost
(Chen \& Guestrin, 2016), siguiendo el esquema de Kim y Pyun (2026):

\begin{itemize}[leftmargin=1.8em]
    \item \textbf{Modelo A (referencia):} solo variables temporales
    (año, mes, tendencia).
    \item \textbf{Modelo B (ampliado):} variables temporales +
    \texttt{english\_ratio}, \texttt{H\_norm}, \texttt{HHI},
    \texttt{repetition\_rate}.
\end{itemize}

La mejora en poder explicativo se evalúa comparando el
\term{coeficiente de determinación ajustado}{수정된 결정계수}{adjusted $R^2$}
y el \term{RMSE}{RMSE}{Root Mean Squared Error}.
Se reporta además la importancia de cada variable (\textit{feature importance})
para identificar cuáles contribuyen más a la explicación del ranking.

% ── V. Resultados ───────────────────────────────────────────────────────────
\section*{V. Resultados}

\subsection*{5.1 Tendencia temporal del ratio inglés (2003--2020)}

\todo{Insertar figura: serie de tiempo del ratio inglés por año con línea de tendencia.}

La prueba de Mann-Kendall sobre la serie anual del ratio inglés
arrojó \todo{estadístico $S$ = X, $p$ = X},
\todo{confirmando / no confirmando} la existencia de una tendencia
creciente estadísticamente significativa en el período analizado.

\begin{table}[H]
\centering
\small
\begin{tabular}{lcc}
\toprule
\textbf{Estadístico} & \textbf{Valor} & \textbf{Interpretación} \\
\midrule
Tau de Kendall ($\tau$) & \todo{X.XX} & \todo{tendencia positiva/negativa} \\
$p$-valor              & \todo{X.XX} & \todo{significativo / no significativo} \\
Pendiente de Sen       & \todo{X.XX} & Cambio promedio anual en ratio inglés \\
\bottomrule
\end{tabular}
\caption{Resultado de la prueba de Mann-Kendall sobre el ratio inglés anual}
\end{table}

\subsection*{5.2 Clasificación Solid--Liquid por año}

\todo{Insertar tabla: año, entropy, H\_norm, HHI, n\_active, clasificación.}

\todo{Insertar figura: gráfico de radar o barras apiladas por año mostrando
la distribución léxica y la clasificación Solid/Liquid.}

Los resultados muestran que \todo{X de 18 años} se clasificaron como Liquid
y \todo{X} como Solid.
Los años clasificados como Liquid se concentran en
\todo{describir período: ej. 2012-2016}, coincidiendo con
\todo{fenómeno contextual, ej. boom del Hallyu internacional}.

\subsection*{5.3 Correlación ranking-variables lingüísticas}

\begin{table}[H]
\centering
\small
\begin{tabular}{lccc}
\toprule
\textbf{Variable} & $r_s$ & $p$-valor & \textbf{Interpretación} \\
\midrule
\texttt{english\_ratio}   & \todo{X.XX} & \todo{X.XX} & \todo{} \\
\texttt{H\_norm}          & \todo{X.XX} & \todo{X.XX} & \todo{} \\
\texttt{HHI}              & \todo{X.XX} & \todo{X.XX} & \todo{} \\
\texttt{repetition\_rate} & \todo{X.XX} & \todo{X.XX} & \todo{} \\
\bottomrule
\end{tabular}
\caption{Correlaciones de Spearman entre variables lingüísticas y ranking}
\end{table}

\subsection*{5.4 Comparación de modelos XGBoost}

\begin{table}[H]
\centering
\small
\begin{tabular}{lccc}
\toprule
\textbf{Modelo} & \textbf{Adjusted $R^2$} & \textbf{RMSE} & \textbf{Variables} \\
\midrule
Modelo A (referencia) & \todo{X.XX} & \todo{X.XX} & Temporales \\
Modelo B (ampliado)   & \todo{X.XX} & \todo{X.XX} & Temporales + lingüísticas \\
\midrule
Diferencia            & \todo{+X.X pp} & \todo{$-$X.X\%} & \\
\bottomrule
\end{tabular}
\caption{Comparación de poder explicativo: Modelo A vs.\ Modelo B}
\end{table}

\todo{Insertar figura: importancia de variables (feature importance) del Modelo B.}

% ── VI. Conclusiones ────────────────────────────────────────────────────────
\section*{VI. Implicaciones y trabajo futuro}

Este estudio presenta un \term{corpus longitudinal}{종단 코퍼스}{longitudinal corpus}
de letras de K-pop 2003--2020 y propone un marco estadístico para analizar
su evolución léxica que combina análisis de tendencia temporal,
clasificación Solid--Liquid basada en entropía y HHI, y modelización
explicativa mediante XGBoost.

Los resultados contribuyen a la estadística aplicada de textos en tres dimensiones:
(1) verificación empírica de la tendencia de aumento del inglés en la
música popular coreana, con prueba de significancia formal;
(2) demostración de que la estructura léxica varía entre años de forma
cuantificable y clasificable; y
(3) evidencia del poder explicativo adicional de variables lingüísticas
sobre la posición en el ranking.

Como trabajo futuro se propone:
incorporar análisis de tópicos (LDA) para definir ejes semánticos formales
al estilo de Moon y Pyun (2025);
ampliar el modelo XGBoost con variables de interacción temporal (gradiente
año a año) siguiendo a Kim y Pyun (2026);
y explorar la relación entre cambios léxicos y contexto sociocultural
(lanzamientos internacionales, eventos políticos).

% ── Referencias ─────────────────────────────────────────────────────────────
\section*{Referencias}

\begin{itemize}[leftmargin=1.2em, label={}, itemsep=0.3em]

\item Chen, T., \& Guestrin, C. (2016). XGBoost: A scalable tree boosting system.
\textit{Proceedings of the 22nd ACM SIGKDD International Conference on Knowledge
Discovery and Data Mining}, 785--794.
\url{https://doi.org/10.1145/2939672.2939785}

\item Jung, S. (2014). K-pop, Indonesian fandom, and social media.
\textit{Transformative Works and Cultures}, 17.
\url{https://doi.org/10.3983/twc.2014.0541}

\item Kendall, M. G. (1975). \textit{Rank Correlation Methods} (4th ed.).
Griffin.

\item Kim, C., \& Pyun, J. (2026).
해외 기상 데이터를 활용한 집중호우 설명 모형의 개선: 대안 데이터 활용의 관점에서.
\todo{[Revista y volumen por confirmar]}.

\item Lee, J. (2004). Linguistic hybridization in K-pop: Discourse of
self-assertion and resistance.
\textit{World Englishes}, 23(3), 429--450.

\item Mann, H. B. (1945). Nonparametric tests against trend.
\textit{Econometrica}, 13(3), 245--259.

\item Moody, A., \& Matsumoto, Y. (2011). Don't touch my moustache: Language
blending and code ambiguation by two J-pop artists.
\textit{Asian Englishes}, 14(1), 4--33.

\item Moon, J., \& Pyun, J. (2025).
리뷰 텍스트 임베딩 기반 소비자 가치축 분류--Entropy·HHI기반 Solid--Liquid 소비유형.
\textit{글로벌경영학회지}, 22(X), 1--1.
\url{https://doi.org/10.38115/asgba.2025.22.x.x}

\item Myers-Scotton, C. (1993). \textit{Duelling Languages: Grammatical Structure
in Codeswitching}. Clarendon Press.

\item Napier, K., \& Shamir, L. (2018). Quantitative sentiment analysis of
lyrics in popular music.
\textit{Journal of Popular Music Studies}, 30(4), 161--176.

\item Schedl, M. (2014). Leveraging microblogs for spatiotemporal music
information retrieval.
\textit{Proceedings of the 36th European Conference on IR Research}, 796--801.

\item Templin, M. (1957). \textit{Certain Language Skills in Children}.
University of Minnesota Press.

\end{itemize}

% ── Glosario ────────────────────────────────────────────────────────────────
\newpage
\section*{Glosario de términos técnicos}

\renewcommand{\arraystretch}{1.2}
\begin{longtable}{p{3.5cm} p{3.5cm} p{3.5cm} p{3cm}}
\toprule
\textbf{한국어} & \textbf{Español} & \textbf{English} & \textbf{Nota} \\
\midrule
\endhead
대중가요 가사 & letras de música popular & popular song lyrics & K-pop en contexto \\
어휘 다양성 & diversidad léxica & lexical diversity & medida con TTR, entropía \\
영어 혼용 & mezcla de inglés & English mixing / code-mixing & tokens en inglés sobre total \\
엔트로피 & entropía & entropy & $H = -\sum p_i \log p_i$ \\
정규화 엔트로피 & entropía normalizada & normalized entropy & $H_{norm} = H / \log K$ \\
허핀달--허쉬만 지수 & índice Herfindahl-Hirschman & HHI & concentración léxica \\
집중형 & tipo concentrado & Solid & un eje léxico dominante \\
분산형 & tipo disperso & Liquid & varios ejes léxicos activos \\
활성 범주 & categoría activa & active category & peso $\geq 10\%$ \\
차트 순위 & posición en el ranking & chart ranking & variable dependiente \\
시계열 분석 & análisis de serie temporal & time series analysis & por año/mes \\
Mann-Kendall 검정 & prueba de Mann-Kendall & Mann-Kendall test & tendencia monotónica \\
스피어만 상관 & correlación de Spearman & Spearman correlation & no paramétrica \\
형태소 분석 & análisis morfológico & morphological analysis & con Kiwipiepy \\
명사 & sustantivo & noun & etiqueta NN* en Kiwi \\
반복률 & tasa de repetición & repetition rate & tokens repetidos / total \\
기준 모형 & modelo de referencia & baseline model & solo variables temporales \\
확장 모형 & modelo ampliado & extended model & + variables lingüísticas \\
변수 중요도 & importancia de variables & feature importance & salida de XGBoost \\
수정된 결정계수 & coef. determinación ajustado & adjusted $R^2$ & penaliza complejidad \\
한류 & ola coreana & Korean Wave / Hallyu & fenómeno cultural global \\
\bottomrule
\end{longtable}

\end{document}
